%==============================================================================
\section{Hồi quy tuyến tính (Linear Regression)}
\label{sec:linear-regression}
%==============================================================================

\begin{dinhnghia}[title={Hồi quy tuyến tính}]
Mô hình thống kê phân tích \textbf{quan hệ tuyến tính} giữa biến phụ thuộc và một hoặc nhiều biến độc lập.
Khi biến độc lập thay đổi, biến phụ thuộc thay đổi tương ứng theo quy luật tuyến tính.
\end{dinhnghia}

\begin{ynghia}[title={Trong học máy}]
Dự đoán giá trị số liên tục từ quan hệ tuyến tính đã học; dùng trong mô hình dự báo, tài chính, đánh giá rủi ro.
\end{ynghia}

\subsection{Dữ liệu ví dụ (một feature)}

\begin{phuongphap}[title={Bảng minh họa}]
Diện tích (ft$^2$, feature $X$) và giá nhà (target $Y$): mô hình dùng $X$ để dự đoán $Y$.
\begin{tabular}{cc}
Diện tích (X) & Giá nhà (Y) \\
1300 & 240 \\
1500 & 320 \\
1700 & 330 \\
1830 & 295 \\
1550 & 256 \\
2350 & 409 \\
1450 & 319 \\
\end{tabular}
\end{phuongphap}

\begin{dinhnghia}[title={Đường hồi quy}]
Đường thẳng thể hiện quan hệ giữa biến phụ thuộc và độc lập.
\end{dinhnghia}

\tsimg{02_linear_regression/img01.jpg}

\subsection{Quan hệ tuyến tính dương và âm}

\begin{ynghia}[title={Quan hệ dương}]
Cả hai biến cùng tăng (hoặc cùng giảm).
\end{ynghia}

\tsimg{02_linear_regression/img02.jpg}

\begin{ynghia}[title={Quan hệ âm}]
Biến độc lập tăng, biến phụ thuộc giảm.
\end{ynghia}

\tsimg{02_linear_regression/img03.jpg}

\subsection{Hai loại hồi quy tuyến tính}

\subsubsection{Hồi quy tuyến tính đơn}

\begin{dinhnghia}[title={Hồi quy tuyến tính đơn}]
Hồi quy tuyến tính đơn là một loại phân tích hồi quy trong đó một biến độc lập duy nhất (còn được gọi là feature) được sử dụng để dự đoán biến phụ thuộc (còn được gọi là target).
\end{dinhnghia}

\tsimg{02_linear_regression/img04.jpg}

\[
\hat{Y} = w_0 + w_1 X + \epsilon
\]

\begin{itemize}
  \item $Y$: target; $X$: feature; $w_0$: hệ số chặn; $w_1$: độ dốc; $\epsilon$: sai số.
\end{itemize}

\subsubsection{Hồi quy tuyến tính đa biến}

\begin{dinhnghia}[title={Hồi quy tuyến tính đa biến}]
Hồi quy tuyến tính đa biến là một loại phân tích hồi quy trong đó nhiều biến độc lập (còn được gọi là features) được sử dụng để dự đoán biến phụ thuộc (còn được gọi là target).
\end{dinhnghia}

\[
\hat{Y} = w_0 + w_1 X_1 + w_2 X_2 + \cdots + w_p X_p + \epsilon
\]

\begin{itemize}
  \item $Y$: target; $X_1, X_2, \ldots, X_p$: features; $w_0, w_1, \ldots, w_p$: hệ số; $\epsilon$: sai số.
\end{itemize}

\subsection{Cách hoạt động}

\begin{dongco}[title={Mục tiêu chính của hồi quy tuyến tính}]
Mục tiêu chính của hồi quy tuyến tính là tìm đường thẳng phù hợp nhất đi qua một tập hợp các điểm dữ liệu sao cho sự khác biệt giữa giá trị thực tế và giá trị dự đoán là nhỏ nhất.
\end{dongco}

\begin{phuongphap}[title={Ba bước chính}]
\begin{enumerate}
  \item \textbf{Giả thuyết}: quan hệ tuyến tính giữa đầu vào và đầu ra.
  \item \textbf{Hàm mất mát}: đo sai số dự đoán (MSE, MAE, \ldots).
  \item \textbf{Tối ưu}: cập nhật tham số đến khi mất mát nhỏ.
\end{enumerate}
\end{phuongphap}

\subsection{Hàm giả thuyết}

\[
\hat{Y} = w_0 + w_1 X \quad \text{(đơn biến)}
\]
\[
\hat{Y} = w_0 + w_1 X_1 + \cdots + w_p X_p \quad \text{(đa biến)}
\]

\subsection{Đường khớp tốt nhất}

Đường hồi quy tốt nhất khi sai số giữa giá trị thực và dự đoán là nhỏ nhất (trung bình).

\tsimg{02_linear_regression/img05.jpg}

\subsection{Hàm mất mát}

\begin{dinhnghia}[title={MSE (phổ biến)}]
\[
J(w_0, w_1) = \frac{1}{2n} \sum_{i=1}^{n} \left( Y_i - \hat{Y}_i \right)^2,
\quad \hat{Y}_i = w_0 + w_1 X_i
\]
\end{dinhnghia}

\subsection{Gradient descent}

\begin{phuongphap}[title={Cập nhật tham số}]
\[
w_0 \leftarrow w_0 - \alpha \frac{\partial J}{\partial w_0},
\quad
w_1 \leftarrow w_1 - \alpha \frac{\partial J}{\partial w_1}
\]
với $\alpha$ là learning rate. Lặp đến khi thay đổi $w_0, w_1$ không đáng kể.
\end{phuongphap}

\subsection{Giả định mô hình}

\begin{luuy}[title={Giả định}]
\begin{itemize}
  \item Ít hoặc không đa cộng tuyến.
  \item Ít tự tương quan trong phần dư.
  \item Quan hệ giữa response và feature là \textbf{tuyến tính}.
\end{itemize}
Vi phạm giả định có thể làm ước lượng chệch hoặc kém hiệu quả.
\end{luuy}

\subsection{Metric đánh giá}

\begin{tinhchat}[title={$R^2$, MSE, RMSE, MAE}]
\[
R^2 = 1 - \frac{\sum (y_i - \hat{y}_i)^2}{\sum (y_i - \bar{y})^2},
\quad
\mathrm{MSE} = \frac{1}{n}\sum_{i=1}^n (y_i - \hat{y}_i)^2
\]
\[
\mathrm{RMSE} = \sqrt{\mathrm{MSE}},
\quad
\mathrm{MAE} = \frac{1}{n}\sum_{i=1}^n |y_i - \hat{y}_i|
\]
\end{tinhchat}

\subsection{Ứng dụng, ưu điểm và thách thức}

\begin{ynghia}[title={Ứng dụng}]
Mô hình dự báo (bất động sản), chọn feature, dự báo tài chính, quản trị rủi ro.
\end{ynghia}

\begin{tinhchat}[title={Ưu điểm}]
Dễ giải thích, huấn luyện nhanh, nền tảng phân tích dự đoán, đơn giản và hiệu quả tính toán.
\end{tinhchat}

\begin{luuy}[title={Thách thức}]
Overfitting, đa cộng tuyến, ảnh hưởng outlier; với dữ liệu phi tuyến có thể cần \textbf{hồi quy đa thức}.
\end{luuy}
